\documentclass[UTF8]{ctexart}
\usepackage{geometry}
\usepackage[backend=biber, style=authoryear]{biblatex} % 示例配置
\addbibresource{references.bib} % 替换为你的参考文献文件
\geometry{margin=1.5cm, vmargin={0pt,1cm}}
\setlength{\topmargin}{-1cm}
\setlength{\paperheight}{29.7cm}
\setlength{\textheight}{25.3cm}

% useful packages.
\usepackage{amsfonts}
\usepackage{amsmath}
\usepackage{amssymb}
\usepackage{amsthm}
\usepackage{enumerate}
\usepackage{graphicx}
\usepackage{multicol}
\usepackage{fancyhdr}
\usepackage{layout}
\usepackage{listings}
\usepackage{xcolor}
%\lstset{language=C++}
\setlength{\headheight}{13pt}  % 增加头部高度
\usepackage{float, caption}

% Listings setting for code formatting
\lstset{
    basicstyle=\ttfamily,
    keywordstyle=\color{blue},
    commentstyle=\color{gray},
    stringstyle=\color{red},
    breaklines=true,
    frame=single,
    numbers=left,
    numberstyle=\tiny,
    stepnumber=1,
    tabsize=4,
}

% some common commands
\newcommand{\dif}{\mathrm{d}}
\newcommand{\avg}[1]{\left\langle #1 \right\rangle}
\newcommand{\difFrac}[2]{\frac{\dif #1}{\dif #2}}
\newcommand{\pdfFrac}[2]{\frac{\partial #1}{\partial #2}}
\newcommand{\OFL}{\mathrm{OFL}}
\newcommand{\UFL}{\mathrm{UFL}}
\newcommand{\fl}{\mathrm{fl}}
\newcommand{\op}{\odot}
\newcommand{\Eabs}{E_{\mathrm{abs}}}
\newcommand{\Erel}{E_{\mathrm{rel}}}

\begin{document}

\pagestyle{fancy}
\fancyhead{}
\lhead{陈梓锐}
\chead{22435036}
\rhead{\today}


%\begin{abstract}
 %   The abstract is not necessary for the theoretical homework, 
   % but for the programming project, 
    %you are encouraged to write one.      
%\end{abstract}

% ============================================

\subsection*{Exercise 11.36}
考虑三阶ERK

$
\left\{
    \begin{aligned}
        y_1 &= f(U^n, t_n) \\
        y_2 &= f(U^n + ka_{21}y_1, t_n + c_2k) \\
        y_3 &= f(U^n + ka_{32}y_2 + ka_{31}y_1, t_n + c_3k) \\
        U^{n+1} &= U^n + k(b_1y_1 + b_2y_2 + b_3y_3) 
    \end{aligned}
\right.
$

$L_u(t_n) = u(t_{n+1}) - u(t_n) - k(b_1y_1 + b_2y_2 + b_3y_3)$

$u(t_{n+1}) = u(t_n) + ku^{'}(t_n) + \frac{k^2}{2}u^{''}(t_n) + \frac{k^3}{6}u^{'''}(t_n) + o(k^3)$

$y_1 = u^{'}(t_n)$

$y_2 = f + kc_2(f_uf + f_t) + \frac{k^2}{2}c_2^2(f_{uu}f^2 + f_{ut}f + f_tf) + o(k^2) $

$y_3 = f + kc_3(f_uf + f_t) + k^2(\frac{c_3^2}{2}(f_{uu}f^2 + f_{ut}f + f_tf) + a_{32}c_2(f^2_uf + f_uf_t)) + o(k^2) $

即\begin{align*}
    L_u(t_n) &= k(1-b_1-b_2-b_3)u^{'} + k^2(\frac{1}{2} - b_2c_2 - b_3c_3)u^{''} \\    
             &= k^3(\frac{1}{6} - \frac{1}{2}b_2c_2^2 - \frac{1}{2}b_3c_3^2)(f_{uu}f^2 + f_{ut}f + f_tf) \\
             &= k^3(\frac{1}{6} - b_3a_{32}c_2)(f_u^2 + f_uf_t) + o(k^3) 
\end{align*}

$\Rightarrow
\left\{
    \begin{aligned}
        b_1 + b_2 + b_3 &= 1 \\
        b_2c_2 + b_3c_3 &= \frac{1}{2} \\
        b_2c_2^2 + b_3c_3^2 &= \frac{1}{3} \\
        b_3a_{32}c_2 &= \frac{1}{6}
    \end{aligned}
\right.
$

将所欲证明的Butcher表代入，发现符合方程要求，即这确实是3-order ERK,
而Heun's方法的$a_{21} = \frac{1}{3}$，并不符合要求。

\subsection*{Exercise 11.42}
"$\Rightarrow$"

$I_s(f) = k\sum_{j=1}^{s}b_jf(t_n+c_jk) = \int_{t_n}^{t_n+k} f(t) \,dx $
对所有r-1阶多项式成立。

归纳法，当r=1时，令$f(x) = 1, I_s(f) = kb_1 = k \Rightarrow b_1 = 1$.
假设RK方法满足$B(i), \forall i = 1,2,\cdots,r-1$,令$f(x) = x^{r-1}, 
I_s(f) = k\sum_{j=1}^{s}b_jf(t_n+c_jk)^{r-1} = \frac{1}{r}((t_n+k)^r - t_n^r)$,二项式展开可得，
\begin{equation}
    \label{eq:11.42}
    k\sum_{j=1}^{s}b_j((c_jk)^{r-1} + \cdots) = \frac{1}{r}(k^r + \cdots).
\end{equation}
下证$\cdots$中的第i项展开即是B(r-1-i)的形式
\begin{align*}
    k\sum_{j=1}^{s}b_jC_{r-1}^i(c_jk)^{r-1-i}\cdot t_n^i &= \frac{1}{r}C_r^i k^{r-i} \cdot t_n^i \\
    \sum_{j=1}^{s}b_j\frac{(r-1)!}{i!(r-1-i)!}c_j^{r-1-i} &= \frac{(r-1)!}{i!(r-i)!} \\
    \sum_{j=1}^{s}b_jc_j^{r-1-i} &= \frac{1}{r-i}
\end{align*}
由式（\ref{eq:11.42}），和归纳假设$\Rightarrow \sum_{j=1}^{s}b_jc_j^{r-1} = \frac{1}{r} \Rightarrow B(r)$.

“$\Leftarrow$” 

只需把必要性的证明倒叙即可。


\subsection*{Exercise 11.59}
在Definition11.56中构造的s阶Hermite插值多项式满足：

$
\left\{
    \begin{aligned}
        P(t_n) &= U^n \\
        P^{'}(t_N^{(i)}) &= f(P(t_n^{(i)}), t_n^{(i)}), \forall i = 1,2,\cdots,s
    \end{aligned}
\right.
$ $t_n^{(i)} = t_n + c_ik$

此时$P(t_{n+1}) = u(t_n) + k\Phi(u(t_n), t_n, k)$
\begin{align*}
    L_u(t_n) &= u(t_{n+1}) - P(t_{n+1}) \\
             &= \frac{1}{(s+1)!}u^{(s+1)}(\xi)(t_{n+1} - t_n)\prod _{i=1}^{s} (t_{n+1} - t_n - c_ik) \\
             &= \frac{1}{(s+1)!}u^{(s+1)}(\xi)k^{s+1}\prod_{i=1}^{s}(1 - c_i)
\end{align*}
这说明s阶配点法至少是s阶精度的。

\subsection*{Exercise 11.60}
因为$\sum_{j=1}^{s}l_j(r) = 1$,

所以$\sum_{j=1}^{s}b_j = \int_{0}^{1} \sum_{j=1}^{s} l_j(r) \,dr  = 1, 
\sum_{j=1}^{s} a_{ij} = \int_{0}^{c_i} \sum_{j=1}^{s} l_j(r) \,dr = c_i$

\subsection*{Exercise 11.62}
\begin{align*}
    l_1(r) &= \frac{(r-\frac{1}{2})(r-\frac{3}{4})}{(\frac{1}{4}-\frac{1}{2})(\frac{1}{4}-\frac{3}{4})} = 8r^2 - 10r + 3 \\
    l_2(r) &= \frac{(r-\frac{1}{4})(r-\frac{3}{4})}{(\frac{1}{2}-\frac{1}{4})(\frac{1}{2}-\frac{3}{4})} = -16r^2 + 16r - 3 \\
    l_3(r) &= \frac{(r-\frac{1}{4})(r-\frac{1}{2})}{(\frac{3}{4}-\frac{1}{4})(\frac{3}{4}-\frac{1}{2})} = 8r^2 - 6r + 1 \\
    a_{11} &= \int_{0}^{\frac{1}{4}} l_1(r) \,dr = \frac{23}{48}, 
    a_{12} = \int_{0}^{\frac{1}{4}} l_2(r) \,dr = -\frac{1}{3}, 
    a_{13} = \int_{0}^{\frac{1}{4}} l_3(r) \,dr = \frac{5}{48} \\
    a_{21} &= \int_{0}^{\frac{1}{2}} l_1(r) \,dr = \frac{7}{12},
    a_{22} = \int_{0}^{\frac{1}{2}} l_2(r) \,dr = -\frac{1}{6},
    a_{23} = \int_{0}^{\frac{1}{2}} l_3(r) \,dr = \frac{1}{12} \\
    a_{31} &= \int_{0}^{\frac{3}{4}} l_1(r) \,dr = \frac{9}{16},
    a_{32} = \int_{0}^{\frac{3}{4}} l_2(r) \,dr = 0,
    a_{33} = \int_{0}^{\frac{3}{4}} l_3(r) \,dr = \frac{3}{16} \\
    b_1 &= \int_{0}^{1} l_1(r) \,dr = \frac{2}{3},
    b_2 = \int_{0}^{1} l_2(r) \,dr = -\frac{1}{3},
    b_3 = \int_{0}^{1} l_3(r) \,dr = \frac{2}{3}
\end{align*}

\begin{table}[H]
    \centering
    \begin{tabular}{c|ccc}
        $\frac{1}{4}$ & $\frac{23}{48}$ & $-\frac{1}{3}$ & $\frac{5}{48}$ \\
        $\frac{1}{2}$ & $\frac{7}{12}$ & $-\frac{1}{6}$ & $\frac{1}{12}$ \\
        $\frac{3}{4}$ & $\frac{9}{16}$ & 0 & $\frac{3}{16}$ \\
        \hline
                    &$\frac{2}{3}$ & $-\frac{1}{3}$ & $\frac{2}{3}$
    \end{tabular}
\end{table}


\subsection*{Exercise 11.65}
依题意，需要用矩阵向量乘法证明B(s+r),C(s)推出D(r)。
\begin{align*}
    (V(c_1,\cdots,c_s)u)_i &= \sum_{k=1}^{s}c_k^{i-1}\sum_{j=1}^{s}b_jc_j^{m-1}a_{jk} \\
                           &= \sum_{k=1}^{s}b_jc_j^{m-1}\sum_{k=1}^{s}c_k^{i-1}a_{jk} \\
                           &= \sum_{j=1}^{s}b_jc_j^{m-1}\frac{c_j^i}{i} \\
                           &= \frac{1}{i}\sum_{j=1}^{s}b_jc_j^{i+m-1}\\ 
                           &= \frac{1}{i(i+m)}
\end{align*}

\begin{align*}
    (V(c_1,\cdots,c_s)v)_i &= \sum_{j=1}^{s}c_j^{i-1}\frac{1}{m}b_j(1-c_j^m) \\
                           &= (\sum_{j=1}^{s}b_jc_j^{i-1} - \sum_{j=1}^{s}b_jc_j^{i+m-1}) \\
                           &= \frac{1}{im} - \frac{1}{(i+m)m} \\
                           &= \frac{1}{i(i+m)}
\end{align*}
即$V(c_1,\cdots,c_s)u = V(c_1,\cdots,c_s)v$，
因为上式只对$\forall i=1,\cdots,s, i+m\leq s$成立，且范德蒙矩阵满秩，
故

$\left\{
    \begin{aligned}
        \forall i& = 1,\cdots,s, \\
        \forall m& = 1,\cdots,r
    \end{aligned}
\sum_{j=1}^{s}b_jc_j^{m-1}a_{ji} = \frac{b_i}{m}(1-c_i^m)
\right.
$


\subsection*{Exercise 11.69}
只需求方法的BCD，使用Theorem11.68，即可得配点法的精度。

由$\sum_{j=1}^{s}b_jc_j^0=1, \sum_{j=1}^{s}b_jc_j^1=\frac{1}{2}, \sum_{j=1}^{s}b_jc_j^2=\frac{1}{3}, 
\sum_{j=1}^{s}b_jc_j^3=\frac{1}{4},\sum_{j=1}^{s}b_jc_j^4 \neq \frac{1}{5}$可知方法是B(4)的。
所以方法精度至多为4.

当m=1时，$\sum_{j=1}^{s}a_{ij} = c_i$,

当m=2时，

$a_{11}c_1 + a_{12}c_2 + a_{13}c_3 = \frac{1}{32}=\frac{c_1^2}{2}$

$a_{21}c_1 + a_{22}c_2 + a_{23}c_3 = \frac{1}{8}=\frac{c_2^2}{2}$

$a_{31}c_1 + a_{32}c_2 + a_{33}c_3 = \frac{9}{32}=\frac{c_3^2}{2}$

当m=3时，

$a_{11}c_1^2 + a_{12}c_2^2 + a_{13}c_3^2 = \frac{1}{192}=\frac{c_1^3}{3}$

$a_{21}c_1^2 + a_{22}c_2^2 + a_{23}c_3^2 = \frac{1}{24}=\frac{c_2^3}{3}$

$a_{31}c_1^2 + a_{32}c_2^2 + a_{33}c_3^2 = \frac{9}{64}=\frac{c_3^3}{3}$

所以方法C(3),由Exercise 11.65可知方法是D(1)的.

由Theorem11.68可知，方法的精度至少为4.










% ===============================================

\end{document}